\documentclass[12pt]{article}

\usepackage{sbc-template}
\usepackage{graphicx,url}
\usepackage{float}
\usepackage[utf8]{inputenc}
\usepackage[brazil]{babel}

\sloppy

\raggedbottom

\setlength{\parskip}{3pt}

\renewcommand{\inst}[1]{}

\AtBeginEnvironment{thebibliography}{%
  \setlength{\itemsep}{0pt plus 1pt minus 0pt}%
  \setlength{\parsep}{0pt}%
}

\title{Fatores associados \`a gravidade de acidentes em rodovias federais
  brasileiras: relev\^ancia preditiva e estabilidade entre macrorregi\~oes}

\author{Diogo Leonel dos Santos (15580980)\\ Gabriel Trigueiro de Melo
  (15511832)\\ Dhener Alves Bento da Silva (15526067)}


\address{Universidade de S\~ao Paulo (USP), S\~ao Paulo, SP, Brasil}

\begin{document}

\maketitle

\begin{resumo}
Modelos de aprendizado de m\'aquina treinados sobre dados abertos da
Pol\'icia Rodovi\'aria Federal (PRF) \cite{prfdados} s\~ao empregados para
identificar fatores de risco associados \`a gravidade de acidentes. Um
estudo recente \cite{gabardo2025} treinou kNN, MLP e Random Forest sobre
registros da regi\~ao Sul (2021--2024), reportou AUC-ROC de 0,99 para a
classe fatal em dois dos tr\^es modelos e derivou fatores de risco via
SHAP \cite{lundberg2017}. Este trabalho quantifica a relev\^ancia
preditiva dos fatores em uma base nacional de registros e avalia a
estabilidade do desempenho e do ordenamento de fatores entre as cinco
macrorregi\~oes por meio de uma matriz de transfer\^encia $5\times 5$. A
concord\^ancia entre os valores SHAP \'e medida pelo tau de Kendall
\cite{kendall1938} e pela sobreposi\c{c}\~ao top-$k$.
\end{resumo}


\section{Introdu\c{c}\~ao}

O tr\^ansito \'e uma das principais causas de morte evit\'avel no Brasil.
Segundo o Anu\'ario Estat\'istico de Sinistros de Tr\^ansito 2024 da PRF
\cite{prf2024}, as rodovias federais registraram 6.160 mortes em 2024, o
maior n\'umero em cinco anos, das quais 4.291 em ocorr\^encias em pista
simples. O Brasil \'e signat\'ario da Agenda 2030 da ONU \cite{onu2015},
cuja meta ODS 3.6 prev\^e reduzir pela metade as mortes e les\~oes no
tr\^ansito at\'e 2030, tema tamb\'em alinhado ao ODS 11. Nesse contexto,
os microdados abertos da PRF \cite{prfdados} tornaram-se insumo frequente
para modelos preditivos que buscam orientar a aloca\c{c}\~ao de recursos
de fiscaliza\c{c}\~ao e de engenharia vi\'aria.

Quando um fator \'e apontado como determinante da letalidade e essa
conclus\~ao orienta investimento p\'ublico, a validade da
recomenda\c{c}\~ao depende da validade do protocolo que a produziu. Esse
protocolo pode estar restrito a um recorte ou enviesado pelo
desbalanceamento dos dados utilizados, levando a predi\c{c}\~oes da
gravidade dos acidentes que podem n\~ao se sustentar em outro contexto.
O estudo de \cite{gabardo2025} ilustra essa situa\c{c}\~ao e,
ao mesmo tempo, \'e exemplar em transpar\^encia, pois os autores
disponibilizaram publicamente o c\'odigo-fonte \cite{gabardo2025codigo}.
O estudo restringe-se \`a regi\~ao Sul, enquanto as recomenda\c{c}\~oes
derivadas dele s\~ao formuladas em termos gerais.

Sob o prisma macrorregional brasileiro, essa limita\c{c}\~ao tem
consequ\^encia direta na formula\c{c}\~ao de pol\'iticas p\'ublicas, e
\'e isso que motiva este trabalho. Frota, malha vi\'aria, clima e
padr\~ao de tr\'afego diferem substancialmente entre as macrorregi\~oes
brasileiras, e n\~ao h\'a evid\^encia estat\'istica de que um conjunto de
fatores estimado no Sul descreva o que acontece no Nordeste ou no Norte.

\textbf{Objetivo.} Identificar, por meio de algoritmos de aprendizado de
m\'aquina, os principais fatores associados \`a gravidade de acidentes em
rodovias federais brasileiras, comparando sua relev\^ancia e capacidade de
predi\c{c}\~ao entre as cinco macrorregi\~oes do IBGE \cite{ibge2017}.
Segundo o crit\'erio de \cite{wazlawick2009}, o objetivo \'e
verific\'avel ao final do trabalho: a relev\^ancia \'e medida por valores
SHAP, a capacidade de predi\c{c}\~ao pelo F1 da classe positiva e a
compara\c{c}\~ao entre regi\~oes pelo tau de Kendall e pela
sobreposi\c{c}\~ao top-$k$ entre ordenamentos.

\section{Trabalhos Relacionados}

\textbf{Predi\c{c}\~ao de gravidade no Brasil.} Diversos trabalhos
aplicam aprendizado de m\'aquina a microdados da PRF, em geral com recorte
geogr\'afico restrito. \cite{mafort2025} modelaram a gravidade de
acidentes em rodovias federais do Rio de Janeiro (2020--2023) comparando
sete classificadores, com regress\~ao log\'istica alcan\c{c}ando 73,85\%
de acur\'acia e interpreta\c{c}\~ao via LIME. \cite{andrade2019}
analisaram a tend\^encia do n\'umero de v\'itimas nas rodovias federais
por macrorregi\~ao antes e depois da D\'ecada de A\c{c}\~ao pela
Seguran\c{c}a no Tr\^ansito e encontraram heterogeneidade regional
consistente ao longo do tempo. O estudo de \cite{gabardo2025}
insere-se nessa linha, restringe-se \`a regi\~ao Sul e aponta como
dire\c{c}\~ao futura a extens\~ao a outras regi\~oes do pa\'is.

\textbf{Generaliza\c{c}\~ao geogr\'afica.} A transferibilidade espacial de
modelos de seguran\c{c}a vi\'aria \'e tema estabelecido
internacionalmente. \cite{man2022} avaliam a portabilidade espacial
e temporal de modelos de previs\~ao de sinistros e recorrem a aprendizado
por transfer\^encia para conter a degrada\c{c}\~ao observada. A origem do
problema \'e a heterogeneidade n\~ao observada: segundo
\cite{mannering2016}, fatores contribuintes variam de forma
sistem\'atica e n\~ao capturada entre contextos, o que gera estimativas
viesadas quando um modelo ajustado em uma regi\~ao \'e aplicado a outra.

\textbf{Transferibilidade de modelos.} Avaliar quanto um modelo ajustado
em um contexto preserva de desempenho em outro constitui linha
metodol\'ogica pr\'opria. \cite{xu2014} tratam o problema por
atualiza\c{c}\~ao bayesiana e mostram que o grau de consist\^encia entre
as distribui\c{c}\~oes de origem e destino condiciona a transferibilidade,
e que a recalibra\c{c}\~ao com amostra reduzida da regi\~ao-alvo recupera
parte do desempenho perdido. \cite{man2022} usam aprendizado por
transfer\^encia com o mesmo fim e reportam degrada\c{c}\~ao
sistem\'atica sob aplica\c{c}\~ao direta e ganhos sob adapta\c{c}\~ao.
Este trabalho difere dessa linha no prop\'osito da medi\c{c}\~ao. Nesses
estudos a degrada\c{c}\~ao \'e um problema a ser mitigado; aqui ela \'e o
objeto de interesse, pois a magnitude da queda entre macrorregi\~oes
indica se recomenda\c{c}\~oes derivadas de um recorte regional podem ser
generalizadas ao pa\'is. A matriz de transfer\^encia completa \'e usada,
portanto, como ferramenta de diagn\'ostico.

\textbf{Posicionamento.} Nenhum dos trabalhos citados mede a relev\^ancia
dos fatores de gravidade sob protocolo corrigido nem avalia a estabilidade
desse ordenamento entre as cinco macrorregi\~oes com controle de tamanho
amostral. Este trabalho pretende preencher essa lacuna.

\section{M\'etodo}

\textbf{Dados.} Utilizam-se os microdados abertos da PRF \cite{prfdados}
na vers\~ao agrupada por pessoa, variante que inclui todas as causas e
tipos de acidente, referentes aos anos de 2021 a 2024, mesmo intervalo
temporal adotado por \cite{gabardo2025}. Os arquivos CSV
(ISO-8859-1) apresentam esquema id\^entico nos quatro anos, com 37
atributos. A base bruta re\'une 2.117.994 registros de envolvidos,
correspondentes a 270.097 acidentes distintos e 652.411 pessoas
distintas. Cada registro \'e atribu\'ido a uma das cinco macrorregi\~oes
do IBGE \cite{ibge2017} pela unidade federativa da ocorr\^encia. Na base
anal\'itica, ap\'os as redu\c{c}\~oes descritas adiante, a
distribui\c{c}\~ao \'e: Sudeste, 30,9\%; Sul, 29,7\%; Nordeste, 21,4\%;
Centro-Oeste, 12,3\%; e Norte, 5,7\%, uma assimetria de mais de cinco
vezes entre a maior e a menor macrorregi\~ao. A vari\'avel resposta \'e
bin\'aria e deriva do atributo de estado f\'isico do envolvido, com valor
positivo quando o desfecho \'e \'obito. Na base anal\'itica, com 608.921
registros, a classe positiva corresponde a 3,72\%, o que d\'a uma
raz\~ao de desbalanceamento de aproximadamente 1 para 26.

\textbf{Algoritmos de aprendizado de m\'aquina.} A etapa de
replica\c{c}\~ao emprega os tr\^es algoritmos usados por
\cite{gabardo2025}: k-vizinhos mais pr\'oximos \cite{cover1967},
perceptron multicamadas \cite{rumelhart1986} e florestas aleat\'orias
\cite{breiman2001}. Essa condi\c{c}\~ao \'e necess\'aria para que a
compara\c{c}\~ao com os valores publicados seja leg\'itima. As an\'alises
seguintes acrescentam \emph{gradient boosting} sobre \'arvores, na
implementa\c{c}\~ao XGBoost \cite{chen2016}, e uma m\'aquina de vetores
de suporte com kernel linear \cite{cortes1995}. O objetivo da escolha
n\~ao \'e maximizar desempenho, mas cobrir fam\'ilias com capacidades de
memoriza\c{c}\~ao distintas. Modelos baseados em dist\^ancia e em margem
linear ajustam-se a fronteiras globais, enquanto conjuntos de \'arvores
particionam o espa\c{c}o de atributos com granularidade suficiente para
isolar observa\c{c}\~oes individuais. Como o ganho esp\'urio produzido por
contamina\c{c}\~ao do conjunto de avalia\c{c}\~ao cresce com a capacidade
de memoriza\c{c}\~ao, comparar as fam\'ilias sob as oito
configura\c{c}\~oes do delineamento permite separar o efeito do protocolo
do efeito do modelo. Os hiperpar\^ametros dos tr\^es algoritmos
replicados s\~ao fixados nos valores reportados por
\cite{gabardo2025}. Para os dois algoritmos acrescentados, s\~ao
definidos por uma busca conduzida uma \'unica vez sobre a
configura\c{c}\~ao integralmente corrigida e mantidos constantes nas
demais, de modo que nenhuma diferen\c{c}a observada entre
configura\c{c}\~oes decorra de reotimiza\c{c}\~ao. Adota-se o kernel
linear porque o custo de treino dos kernels n\~ao lineares cresce de
forma superquadr\'atica com o n\'umero de amostras, o que \'e
incompat\'ivel com as 33 execu\c{c}\~oes previstas entre o delineamento
fatorial e a matriz de transfer\^encia.

\textbf{Prepara\c{c}\~ao e decis\~oes de escopo.} Tr\^es decis\~oes
definem a base anal\'itica. Primeiro, a variante replica cada envolvido
por causa e por tipo: a base bruta tem 2.117.994 linhas para 652.411
pessoas; a restri\c{c}\~ao \`a causa principal reduz esse total a
1.228.893, e a restri\c{c}\~ao ao tipo prim\'ario resulta em 716.779
registros, uma linha por envolvido por ocorr\^encia. A multiplicidade
n\~ao compromete a independ\^encia entre parti\c{c}\~oes, j\'a que todas
as linhas de um mesmo envolvido compartilham o identificador do acidente,
mas distorce a pondera\c{c}\~ao, pois ocorr\^encias mais detalhadas
contribuem desproporcionalmente para o ajuste. Segundo, 15,95\% dos
registros n\~ao possuem r\'otulo utiliz\'avel (valores ausentes e ``n\~ao
informado'') e s\~ao exclu\'idos, pois atribu\'i-los \`a classe negativa
inflaria a maioria; a base anal\'itica re\'une 608.921 registros.
Terceiro, os atributos de contagem de ilesos, feridos leves, feridos
graves e \'obitos s\~ao codifica\c{c}\~oes determin\'isticas do estado
f\'isico (coincid\^encia total com o r\'otulo) e constituem o alvo sob
outra forma. Por isso s\~ao removidos, assim como os identificadores de
registro, pessoa e ve\'iculo. Verificar se esses atributos estavam entre
os preditores de \cite{gabardo2025} faz parte da
replica\c{c}\~ao. Em raz\~ao da estrutura de agrupamento, adota-se
particionamento por grupo pelo identificador do acidente em todas as
oito configura\c{c}\~oes, como decis\~ao de protocolo. A
verifica\c{c}\~ao de integridade confirmou que n\~ao h\'a colis\~ao de
identificadores entre os anos.

\textbf{Implementa\c{c}\~ao.} A an\'alise \'e reimplementada em Python,
com scikit-learn \cite{pedregosa2011} e imbalanced-learn
\cite{lemaitre2017}. Antes de qualquer corre\c{c}\~ao, a
configura\c{c}\~ao original \cite{gabardo2025codigo} \'e reproduzida
integralmente, incluindo seus defeitos, e comparada aos valores
publicados em \cite{gabardo2025}. Sem essa equival\^encia, qualquer
diferen\c{c}a posterior seria amb\'igua entre efeito da corre\c{c}\~ao e
efeito da reimplementa\c{c}\~ao. Considera-se a replica\c{c}\~ao
bem-sucedida quando o F1 da classe positiva difere em no m\'aximo 0,02 do
valor publicado. Registra-se que o \emph{script} disponibilizado pelos
autores \cite{gabardo2025codigo} n\~ao executa sem a corre\c{c}\~ao de um
erro de digita\c{c}\~ao em uma chamada de fun\c{c}\~ao. A etapa de
replica\c{c}\~ao \'e conduzida sobre a variante padr\~ao utilizada pelos
autores, e a diverg\^encia de base fica restrita \`as an\'alises
seguintes.

\textbf{Hiperpar\^ametros e m\'etricas.} Os hiperpar\^ametros s\~ao
fixados nos valores de \cite{gabardo2025}, sem reotimiza\c{c}\~ao
entre configura\c{c}\~oes, porque o objeto de compara\c{c}\~ao \'e o
protocolo e n\~ao o modelo. A m\'etrica prim\'aria \'e o F1 da classe
positiva. Precis\~ao, revoca\c{c}\~ao e \'area sob a curva
precis\~ao-revoca\c{c}\~ao \cite{saito2015} s\~ao reportadas
separadamente, como m\'edia e desvio-padr\~ao entre as dobras de uma
valida\c{c}\~ao cruzada estratificada de cinco dobras. A dispers\~ao
entre dobras tamb\'em \'e informativa: configura\c{c}\~oes contaminadas
tendem a apresentar vari\^ancia anormalmente baixa, o que constitui
evid\^encia de vazamento \cite{kaufman2012,kapoor2023} independente do
valor absoluto da m\'etrica.

\textbf{Generaliza\c{c}\~ao regional.} Constr\'oi-se uma matriz de
transfer\^encia $5\times 5$: para cada macrorregi\~ao $i$, treina-se um
modelo exclusivamente com seus registros e avalia-se sobre cada
macrorregi\~ao $j$, incluindo a pr\'opria. O elemento $M_{ij}$ \'e o F1
da classe positiva. A diagonal representa o desempenho intrarregional
sobre parti\c{c}\~ao retida, e os elementos fora da diagonal representam
a transfer\^encia, com degrada\c{c}\~ao relativa definida como
$(M_{ii}-M_{ij})/M_{ii}$. A matriz completa \'e prefer\'ivel \`a
avalia\c{c}\~ao de uma \'unica origem porque permite distinguir a
hip\'otese de que uma regi\~ao espec\'ifica n\~ao exporta seu modelo da
hip\'otese de que nenhuma regi\~ao exporta; apenas a segunda sustenta a
afirma\c{c}\~ao de que um modelo nacional \'unico \'e inadequado. Como o
volume de registros difere em fator superior a cinco entre a maior e a
menor macrorregi\~ao, todos os conjuntos de treino s\~ao subamostrados
para o $n$ da regi\~ao Norte, preservando a propor\c{c}\~ao de classes,
com repeti\c{c}\~ao em amostragens independentes e agrega\c{c}\~ao dos
resultados. Sem esse controle, a diferen\c{c}a de tamanho amostral
poderia ser interpretada como heterogeneidade regional.

\textbf{Estabilidade dos fatores.} Os valores SHAP \cite{lundberg2017}
s\~ao calculados sobre as parti\c{c}\~oes de valida\c{c}\~ao e
convertidos em ordenamento de import\^ancia. Realizam-se duas
compara\c{c}\~oes: o ordenamento sob protocolo original \emph{versus}
corrigido, na mesma regi\~ao, e os ordenamentos entre macrorregi\~oes sob
protocolo corrigido. A concord\^ancia \'e quantificada pelo tau de
Kendall \cite{kendall1938} entre ordenamentos completos e pela
sobreposi\c{c}\~ao relativa entre os dez atributos mais influentes, com
$k = 5$ como verifica\c{c}\~ao de robustez. Concord\^ancia alta indica
fatores robustos; concord\^ancia baixa indica que o ordenamento \'e
artefato de protocolo ou espec\'ifico de uma regi\~ao.

\textbf{Reprodutibilidade.} Sementes aleat\'orias s\~ao fixadas e
reportadas; c\'odigo e vers\~oes de bibliotecas s\~ao depositados em
reposit\'orio p\'ublico. Os resultados s\~ao consolidados em tabelas de
replica\c{c}\~ao, de decomposi\c{c}\~ao por configura\c{c}\~ao e de
desempenho por modelo, al\'em da matriz de transfer\^encia regional e do
gr\'afico comparativo de estabilidade dos fatores.

\section{Cronograma}

Os marcos, vinculados \`as datas da disciplina em 2026, s\~ao
apresentados na Tabela~\ref{tab:cronograma}.

\begin{table}[H]
\centering
\caption{Cronograma de marcos e entregas}
\label{tab:cronograma}
\begin{tabular}{|p{3.2cm}|p{9.6cm}|}
\hline
\textbf{Prazo} & \textbf{Marco / Atividade} \\ \hline
At\'e 26/08/2026 & Entrega do plano de trabalho. \\ \hline
At\'e 16/09/2026 & Replica\c{c}\~ao da configura\c{c}\~ao original,
primeira medida sob protocolo corrigido e auditoria do conjunto de
preditores do estudo de refer\^encia. \\ \hline
At\'e 21/10/2026 & Relat\'orio parcial e v\'ideo de 8 a 10 minutos:
decomposi\c{c}\~ao fatorial completa e primeiras macrorregi\~oes da
matriz de transfer\^encia. \\ \hline
At\'e 18/11/2026 & Relat\'orio final: matriz de transfer\^encia completa,
compara\c{c}\~ao de estabilidade dos fatores e desempenho consolidado dos
modelos. \\ \hline
19/11/2026 & Entrega dos slides e apresenta\c{c}\~ao presencial de 12
minutos, seguida de 3 minutos de perguntas. \\ \hline
\end{tabular}
\end{table}

\bibliographystyle{sbc}
\bibliography{referencias}

\end{document}
%%% Local Variables:
%%% mode: LaTeX
%%% TeX-master: t
%%% End:
